\section{Thai}

% NOTE the \thaiinline{} around the Thai part of the frame title.
% Headings are set in Fira Sans, which has no Thai glyphs, so unwrapped Thai
% in a title comes out as empty boxes. See preamble/lang-thai.tex.
\begin{frame}{Thai text / \thaiinline{ข้อความภาษาไทย}}
  \begin{block}{Which font are you seeing?}
    Built with \textbf{XeLaTeX} or \textbf{LuaLaTeX}, the Thai below is
    \hl{Sarabun} (TH Sarabun New).\\
    Built with \textbf{pdfLaTeX}, it falls back to \hl{Garuda}.
  \end{block}

  \begin{thaipar}
    สวัสดีครับ นี่คือตัวอย่างข้อความภาษาไทยบนสไลด์
    เพื่อทดสอบว่าฟอนต์แสดงผลถูกต้องหรือไม่
  \end{thaipar}

  \vspace{0.6em}
  Mixed in a sentence: the Thai word \thaiinline{ตัวอย่าง} means ``example''.
\end{frame}

\begin{frame}{Thai list / \thaiinline{รายการภาษาไทย}}
  \begin{itemize}
    \item \thaiinline{ข้อแรก} --- first item
    \item \thaiinline{ข้อที่สอง} --- second item
    \item \thaiinline{ข้อที่สาม} --- third item
  \end{itemize}

  \vspace{0.6em}
  \begin{thaipar}
    การจัดวางตัวอักษรไทยต้องรองรับสระบนและวรรณยุกต์
    เช่น ที่ ปี้ กื๊อ ให้แสดงผลได้ถูกต้อง
  \end{thaipar}

  \vspace{0.4em}
  {\scriptsize\color{BrandMuted}The line above deliberately stacks vowel and
  tone marks --- a good check that the font is rendering properly.}
\end{frame}
